% TOP_PROBLEM: {"id": 30007425, "title": "Diagonal extension from positively curved surfaces in R3", "category_id": 8, "category": {"id": 8, "name": "analysis", "display_name": "Analysis", "description": "Limits, continuity, calculus, and function theory.", "slug": "analysis", "order_index": 8, "created_at": "2026-07-31T15:26:25.671Z"}, "status": "open"}
% FIRST_POSTED: 2026-10-06
% ATTEMPT_STATUS: solved
% OPENAI_CLAIM: {"schema_version":1,"source_repo":"https://github.com/openai/math","source_commit":"adc7f1241b42e322a6451854ab7e4b4c146bf78a","release_date":"2026-10-06","match":"full","resolution":"proved","independently_reviewed":false,"second_pass":"agreed","families":[{"family":"077","title":"Fourier restriction for positively curved surfaces","subject":"Real and complex analysis","manuscripts":[{"title":"Diagonal Fourier extension for positively curved surfaces in three dimensions","pdf_path":"preprints/Diagonal-Fourier-extension-for-positively-curved-surfaces-in-three-dimensions-September-24-2026/Diagonal-Fourier-extension-for-positively-curved-surfaces-in-three-dimensions-September-24-2026.pdf","readme_path":"preprints/Diagonal-Fourier-extension-for-positively-curved-surfaces-in-three-dimensions-September-24-2026/README.md","date":"2026-09-24","theorem_ref":"thm:main (sections/00-introduction.tex:18)","statement_sources":[{"path":"preprints/Diagonal-Fourier-extension-for-positively-curved-surfaces-in-three-dimensions-September-24-2026/build/sections/00-introduction.tex","line":18,"label":"thm:main"}]}],"lean":null,"reasoning_trace":null}]}
\documentclass[11pt]{article}
\usepackage[margin=25mm]{geometry}
\usepackage{fontspec}
\setmainfont{Latin Modern Roman}
\usepackage{amsmath,amssymb,mathtools}
\usepackage{unicode-math}
\setmathfont{Latin Modern Math}
\usepackage{enumitem}
\usepackage{xurl,hyperref,bookmark}
\hypersetup{colorlinks=true,urlcolor=blue}
\setcounter{secnumdepth}{0}
\setlength{\parindent}{0pt}
\setlength{\parskip}{5pt}
\setlength{\emergencystretch}{3em}
\begin{document}
\section*{Diagonal extension from positively curved surfaces in R3}

\subsection{OpenAI's claim}
\textit{Extension from curved surfaces}: For a compact smooth surface in $\mathbb R^3$ with positive curvature, the paper claims diagonal Fourier extension boundedness for every exponent $p>3$. It includes smooth surface boundaries and draws a two-dimensional Schrödinger smoothing consequence from the paraboloid case.

\subsection{Mathematical statement of the claimed result}
\paragraph{Mathematical claim 1.} Let $\Sigma\subset\mathbb R^3$ be a compact smooth surface with positive curvature, allowing smooth boundary, and use surface area $\sigma$. For $E_\Sigma f(x)=\int_\Sigma e^{2\pi ix\cdot\xi}f(\xi)\,d\sigma(\xi)$, the claim is \[\|E_\Sigma f\|_{L^p(\mathbb R^3)}\le C_{\Sigma,p}\|f\|_{L^p(\Sigma,\sigma)}\quad(3<p<\infty),\] for all complex-valued $f\in L^p(\Sigma)$. It concerns this diagonal range in three dimensions, with a surface-dependent constant.

\subsection{Outline of the argument}
\textit{Extension from curved surfaces}: Graph patches are extended to phases compatible with an elliptic-capacity propagation theorem. That theorem is combined with weighted Kakeya control to sum packet coefficients by their capacity. A joint estimate for widely separated balls removes the local spatial loss, and chart assembly transfers the estimate to the original surface.

\subsection{Scope relative to this problem}
\textit{Extension from curved surfaces}: The primary statement says: For every compact smooth positively curved surface in R3, allowing smooth boundary, its extension operator maps Lp(surface) to Lp(R3) for every p>3. The target retains these quantified objects and hypotheses, so its positive assertion is covered.

With the hypotheses and conventions above, the claimed result addresses the whole mathematical assertion on this page.

\subsection{Formal verification}
\textit{Extension from curved surfaces}: The pinned catalogue lists no Lean formalization for this result.

This site has not run the Lean development. These links describe OpenAI’s selected formal statements and should be read with their stated scope.

\subsection{Status caveat}
OpenAI reports that an internal model produced the manuscripts. Their arguments have not been independently verified as proofs by this site. Any portion outside the documented formal scope may contain mistakes, and the catalogue label reports the announced mathematical claim.

\subsection{Sources}
\begin{itemize}
\item OpenAI, \emph{Diagonal Fourier extension for positively curved surfaces in three dimensions}, 2026-09-24; citation key \texttt{OAI:Diagonal-Fourier-extension-for-positively-curved-surfaces-in-three-dimensions-September-24-2026}. \href{https://github.com/openai/math/blob/adc7f1241b42e322a6451854ab7e4b4c146bf78a/preprints/Diagonal-Fourier-extension-for-positively-curved-surfaces-in-three-dimensions-September-24-2026/Diagonal-Fourier-extension-for-positively-curved-surfaces-in-three-dimensions-September-24-2026.pdf}{PDF (GitHub)}; \href{https://github.com/openai/math/raw/adc7f1241b42e322a6451854ab7e4b4c146bf78a/preprints/Diagonal-Fourier-extension-for-positively-curved-surfaces-in-three-dimensions-September-24-2026/Diagonal-Fourier-extension-for-positively-curved-surfaces-in-three-dimensions-September-24-2026.pdf}{PDF (direct)}; \href{https://github.com/openai/math/blob/adc7f1241b42e322a6451854ab7e4b4c146bf78a/preprints/Diagonal-Fourier-extension-for-positively-curved-surfaces-in-three-dimensions-September-24-2026/README.md}{Manuscript page and citation}; \href{https://github.com/openai/math/blob/main/preprints/Diagonal-Fourier-extension-for-positively-curved-surfaces-in-three-dimensions-September-24-2026/Diagonal-Fourier-extension-for-positively-curved-surfaces-in-three-dimensions-September-24-2026.pdf}{Latest version}; Statement source: \href{https://github.com/openai/math/blob/adc7f1241b42e322a6451854ab7e4b4c146bf78a/preprints/Diagonal-Fourier-extension-for-positively-curved-surfaces-in-three-dimensions-September-24-2026/build/sections/00-introduction.tex\#L18}{thm:main}.
\item \href{https://github.com/openai/math/blob/adc7f1241b42e322a6451854ab7e4b4c146bf78a/CONTENTS.md}{OpenAI catalogue entry}.
\end{itemize}
\end{document}
